\section{Introduction}

An LLM agent reads from many places. It reads retrieved documents and the output of the tools it calls. It reads the descriptions of those tools, the files and skills it loads, and messages
from other agents. It also reads its own stored memory~\cite{chu-2026-systematic,ma-2025-safety}. All of this arrives in one context window, and the model treats it as one body of
text. The model has no reliable way to tell which part of that text is trustworthy and which part is not~\cite{greshake-2023-not}.

It helps to picture the context window as a body of water fed by a fixed set
of inflow streams. Whatever enters through any stream mixes into the same body
and affects everything downstream, whether it was dumped in deliberately or
arrived as runoff. The question this paper answers first is simply, how
many streams are there, and how many consequences may arise because of this pollution? 

Two research communities study what goes wrong when untrustworthy content gets
in, and they use different words for it. The security community studies content that an attacker placed there on purpose. It names attacks such as indirect prompt
injection~\cite{greshake-2023-not}, retrieval
poisoning~\cite{zou-2024-poisonedrag,zhong-2023-poisoning} and memory
backdoors~\cite{dash-2026-untrusted}, and it publishes defenses against
them~\cite{xiang-2024-certifiably,liu-2025-datasentinel,an-2025-ipiguard}.
Success means the attacker fails to achieve a goal.

The AI/ML community studies content that became wrong without any
attacker. It names mechanisms such as lost-in-the-middle~\cite{liu-2023-lost} and
context rot~\cite{xia-2026-diagnosing}. It also names knowledge conflict, where a retrieved passage
disagrees with what the model already
learned~\cite{xie-2023-adaptive}. It publishes fixes for
them~\cite{huang-2025-parammute,bi-2025-parameters,yuan-2024-discerning}.
Success means accuracy goes up.

The two communities describe the same underlying failure: a passage the model
should not trust enters the context window and changes the output. Only the
provenance differs; the model sees tokens either way. The difference is not merely one of vocabulary. If it were, a reader could
translate between the two communities and lose nothing. They differ in
something load-bearing: what counts as a successful evaluation? The AI/ML community measures a fix for the accidental case on an average-case
benchmark, because nothing adapts to it. The security community measures a defense for the
adversarial case against the worst case, because something does. Those are not two styles
of the same test. A method can pass the first test convincingly and fail the second one
completely. Neither community's publishing process catches this, because each
reviews against its own convention.

So the field publishes, adopts, and cites a method without anyone asking
whether it survives the other community's test. That is a testable claim about the literature, not a complaint about it. We
test it by reading both literatures under one taxonomy and recovering every
experiment either side has run. The test finds a gap. We then run the
experiments that fall in it to answer the following research questions.

\paragraph{Research questions.}
\begin{description}
\item[RQ1] \textbf{Are the two communities reinventing the wheel?} We read
  \NCorpus{} papers from both sides under one scheme, compare where each
  community's defenses intervene, and recover every (defense, attack) pair
  the \NAnalysed{} papers that name a technique report. They are not, and
  what they do instead is what leaves the region untested
  (Sections~\ref{sec:literatures}, \ref{sec:parallel}).
\item[RQ2] \textbf{Does a defense built for one regime's problem survive
  the other's?} We execute the pairs the literature skipped, from the
  authors' released code and weights
  (Sections~\ref{sec:experiments}, \ref{sec:results}).
\item[RQ3] \textbf{Is failure a property of individual methods or of a
  design choice they share?} We build a payload that separates the
  explanations rather than maximizes damage
  (Section~\ref{sec:distractor}).
\item[RQ4] \textbf{Do the measurements the field leaves out change the
  result?} We re-read every result against three arms almost no paper
  reports: what the model scores with no document at all, whether a defense
  gained accuracy or only stopped answering, and whether the result holds on
  a second base model. We also run two implementations of the one principle
  that held, to tell a property of a design from a quirk of one codebase
  (Section~\ref{sec:results2}).
\item[RQ5] \textbf{Does that explanation survive an adaptive attacker?} We
  build two optimizers against the one method that resisted
  (Section~\ref{sec:adaptive}).
\end{description}

Both literatures have been systematized before: the security side
repeatedly~\cite{dehghantanha-2026-sok,chu-2026-systematic,jiang-2026-sok},
and the AI/ML side on its own terms as a problem of knowledge
conflict~\cite{xu-2024-knowledge}. None of it runs the experiment. A
survey of one literature has only that literature's methods in hand, and the
test we need takes a defense from one side and an attack from the other. So we
read both, found the pairs nobody has run, and ran them from the authors'
released code and weights. To our knowledge this is the first evaluation
in which published defenses from one of these communities meet published
failure mechanisms from the other.

\paragraph{Scope.}
Three limits bound what follows. Each setting runs on one dataset, and we plant
poisoned passages into the context rather than attack the retriever, so the
numbers describe what a generator does with bad text. Most runs use one model
family at one size, and our own control shows that a different model can reverse
an outcome. Almost every attack we run knows nothing of the defense it meets,
which bounds attack strength from below: conservative evidence when a defense
fails, and no evidence at all when one holds.

\paragraph{Contributions.}
\begin{itemize}
\item \textbf{We show that the two communities divide the problem rather than
  duplicate it (RQ1).} One side studies how the bad text arrives and defends
  where it enters; the other studies why the model believes it and defends
  where it reasons. That division, not duplication, is what leaves a region
  neither side has tested: of every confirmed (defense, attack) pair either
  community reports, only a handful cross between them, and a security
  defense has been evaluated against an AI/ML failure mechanism exactly
  \NSecDefVsSafetyAttack{} time (Sections~\ref{sec:parallel},
  \ref{sec:emptycell}).

\item \textbf{We build a platform and run \NCampaignPairs{} pairs on it
  (RQ2).} Every method executes against an undefended arm on the same items,
  under controls we fixed before any run (Section~\ref{sec:experiments}).

\item \textbf{We trace the failures to a single design choice (RQ3).}
  Every method that fails weights the retrieved text against the model's
  memory, and turning that weight relocates correctness instead of creating
  it. What they take from memory the attacker receives,
  and what they take from the attacker memory receives, so neither direction
  moves anything toward the real document (Sections~\ref{sec:results},
  \ref{sec:distractor}).

\item \textbf{We re-read every result against three arms the field
  omits, and one finding does not survive them (RQ4).} Without a
  no-retrieval arm a defense that never uses the document looks robust.
  Without an abstention axis a method that stopped answering looks inert.
  Without a second model an effect that reverses looks settled
  (Section~\ref{sec:results2}).

\item \textbf{We build two adaptive attacks against the one defense that
  resisted (RQ5).} It loses \PctAcdRemovedGradient{}\% of its
  protection to an attacker who reads its rule, which prices the method rather
  than dismissing it (Section~\ref{sec:adaptive}).

\item \textbf{We state what a deployment inherits from these methods.} A
  method that tunes how much the model trusts retrieved text is correct when
  the text is merely stale and dangerous when an attacker owns it, and a
  writable corpus is the second case (Section~\ref{sec:discussion}).

\item \textbf{We derive the design rule a defense covering both regimes
  would have to follow, and we have not yet built one.} Every signal inside
  the context window belongs to whoever wrote the passage, which is why the
  reasoning stage fails and why the one method that resisted fell to an
  attacker who read its rule. Such a method has to key on something the
  attacker does not hold --- provenance, agreement across independently
  retrieved passages, or the decision to act rather than the decision to
  believe --- and it still needs the reasoning stage for the stale document
  nobody wrote as an attack. The experiment that tests this is specified and
  queued; it has not run, and we claim no result from it
  (Section~\ref{sec:discussion}).

\item \textbf{We release a runnable artifact.} Every attack and defense
  executes from its authors' released code at a pinned commit, verified
  unmodified; the harness, the payloads and the full census ship with it.
\end{itemize}

\paragraph{How the paper reads.} Section~\ref{sec:background} fixes the
terms and the threat model. Section~\ref{sec:literatures} answers RQ1 and
ends at the region nobody has tested. Section~\ref{sec:experiments} builds
the harness and Section~\ref{sec:results} runs it, which answers RQ2 and,
in Section~\ref{sec:distractor}, RQ3. Section~\ref{sec:results2} tries to
break those results with the controls the literature omits (RQ4) and then
with an adaptive attacker (RQ5). Section~\ref{sec:discussion} turns what
survives into conventions, into what a deployment inherits, and into the
design rule a defense covering both regimes would have to follow.
Section~\ref{sec:threats} states the limits and what is not finished.
